% --- Archon physics formalization source begin ---
% archon:physics
% archon:covers PhyXMiniProblems/problem_phyx_mini_0916.lean
% archon:source-report reports/phyx_mini/problem_phyx_mini_0916.source.json

\chapter{Physics problem phyx\_mini\_0916}
\label{ch:PhyXMiniProblems_problem_phyx_mini_0916}

\paragraph{Problem source.}
\#\# Physical scenario

A copper wire carries a steady 125 \$A\$ current to an electroplating tank.

\#\# Question

Find the magnetic field due to a 1.0 cm segment of this wire at a point 1.2 m away from it, if the point is point \$P\_2\$, in the \$xy\$-plane and on a line at \$30\textasciicircum{}\textbackslash{}circ\$ to the segment.

\#\# Answer choices

- A: A: 4.3 \textbackslash{}times 10\textasciicircum{}\{-8\}

- B: B: -2.15 \textbackslash{}times 10\textasciicircum{}\{-8\}

- C: C: -4.3 \textbackslash{}times 10\textasciicircum{}\{-8\}

- D: D: -4.7 \textbackslash{}times 10\textasciicircum{}\{-8\}

\#\# Figure caption (auxiliary; use the image as primary evidence)

The image depicts a diagram involving a 3D coordinate system with axes labeled \textbackslash{}(x\textbackslash{}), \textbackslash{}(y\textbackslash{}), and \textbackslash{}(z\textbackslash{}). 

- A horizontal wire along the \textbackslash{}(x\textbackslash{})-axis carries a current of \textbackslash{}(125 \textbackslash{}, \textbackslash{}text\{A\}\textbackslash{}) to the left, indicated by an arrow and a label.
- The wire bends around an obstacle, depicted as a gray box on the left side.
- Two points, \textbackslash{}(P\_1\textbackslash{}) and \textbackslash{}(P\_2\textbackslash{}), are marked in space:
  - \textbackslash{}(P\_1\textbackslash{}) is directly above the wire on the \textbackslash{}(y\textbackslash{})-axis, at a distance of \textbackslash{}(1.2 \textbackslash{}, \textbackslash{}text\{m\}\textbackslash{}).
  - \textbackslash{}(P\_2\textbackslash{}) is located at an angle of \textbackslash{}(30\textasciicircum{}\textbackslash{}circ\textbackslash{}) from the \textbackslash{}(x\textbackslash{})-axis, at the same vertical distance of \textbackslash{}(1.2 \textbackslash{}, \textbackslash{}text\{m\}\textbackslash{}) but diagonally positioned away from the wire.
- \textbackslash{}(P\_2\textbackslash{}) is labeled with a distance of \textbackslash{}(1.2 \textbackslash{}, \textbackslash{}text\{m\}\textbackslash{}).
- A perpendicular distance of \textbackslash{}(1.0 \textbackslash{}, \textbackslash{}text\{cm\}\textbackslash{}) from the \textbackslash{}(z\textbackslash{})-axis to the wire is marked.

The diagram illustrates the spatial relationship between the current-carrying wire and the specified points.

\#\# Dataset metadata

Magnetism; Physical Model Grounding Reasoning, Spatial Relation Reasoning

\paragraph{Recorded answer/context.}
C: C: -4.3 \textbackslash{}times 10\textasciicircum{}\{-8\}

\paragraph{Figure/image path.}
/root/proposal\_for\_physic/science-mango/phyx\_mini\_run/phyx\_data/test\_image/916.png

\paragraph{Formalization target.}
create a compiling Lean file with sorry bodies at `PhyXMiniProblems/problem\_phyx\_mini\_0916.lean`. The Lean declarations must preserve the physical quantities, dimensions or dimensional roles, figure labels, governing-law hypotheses, and final relation expressed by this problem.
Use Mathlib/Physlib names found through LeanExplore where available. If a physics API is missing, introduce faithful local abstractions rather than scalar placeholder aliases.

\begin{theorem}[Physics formalization target]
\label{thm:physics:phyx_mini_0916:target}
\lean{PhyXMiniProblems.ProblemPhyXMini0916.magneticFieldDueToSegmentAtP2}
\uses{def:physics:phyx-mini-0916:phyxminiproblems-problemphyxmini0916-shortwiremagneticfieldsetup, def:physics:phyx-mini-0916:phyxminiproblems-problemphyxmini0916-magneticfieldzcomponentatp2inteslas, def:physics:phyx-mini-0916:phyxminiproblems-problemphyxmini0916-matchesshortwireproblemdata, def:physics:phyx-mini-0916:phyxminiproblems-problemphyxmini0916-matchesprimaryshortwirefigure, def:physics:phyx-mini-0916:phyxminiproblems-problemphyxmini0916-hasphysicalshortwireparameters, def:physics:phyx-mini-0916:phyxminiproblems-problemphyxmini0916-usesstandardvacuumpermeability, def:physics:phyx-mini-0916:phyxminiproblems-problemphyxmini0916-satisfiessteadycurrentfieldmodel, def:physics:phyx-mini-0916:phyxminiproblems-problemphyxmini0916-satisfiesshortsegmentbiotsavartlaw, def:physics:phyx-mini-0916:phyxminiproblems-problemphyxmini0916-displayedmagneticfieldinteslas, def:physics:phyx-mini-0916:phyxminiproblems-problemphyxmini0916-roundstonearesttenthtimestenpownegeight, def:physics:phyx-mini-0916:phyxminiproblems-problemphyxmini0916-isuniqueclosestdisplayedanswer}
The assigned autoformalize agent should translate this physics problem into Lean declarations in the covered file, with theorem and lemma proof bodies written as `by sorry`.
\end{theorem}
\begin{proof}
This is an autoformalization task, not a proof task. Produce statements that can later be proved by the physics prover without weakening the physical model.
\end{proof}

% --- Archon declaration topology begin ---
\section{Declaration topology}
\begin{definition}[electric Current Dimension]
  \label{def:physics:phyx-mini-0916:phyxminiproblems-problemphyxmini0916-electriccurrentdimension}
  \lean{PhyXMiniProblems.ProblemPhyXMini0916.electricCurrentDimension}
  Electric current has dimension charge per unit time
\end{definition}
\begin{proof}
  This is the declared model definition of electric Current Dimension.
\end{proof}
\begin{definition}[magnetic Permeability Dimension]
  \label{def:physics:phyx-mini-0916:phyxminiproblems-problemphyxmini0916-magneticpermeabilitydimension}
  \lean{PhyXMiniProblems.ProblemPhyXMini0916.magneticPermeabilityDimension}
  Magnetic permeability has dimension M L C⁻²
\end{definition}
\begin{proof}
  This is the declared model definition of magnetic Permeability Dimension.
\end{proof}
\begin{definition}[Length Quantity]
  \label{def:physics:phyx-mini-0916:phyxminiproblems-problemphyxmini0916-lengthquantity}
  \lean{PhyXMiniProblems.ProblemPhyXMini0916.LengthQuantity}
  A nonnegative, unit-independent physical length
\end{definition}
\begin{proof}
  This is the declared model definition of Length Quantity.
\end{proof}
\begin{definition}[Electric Current Quantity]
  \label{def:physics:phyx-mini-0916:phyxminiproblems-problemphyxmini0916-electriccurrentquantity}
  \lean{PhyXMiniProblems.ProblemPhyXMini0916.ElectricCurrentQuantity}
  \uses{def:physics:phyx-mini-0916:phyxminiproblems-problemphyxmini0916-electriccurrentdimension}
  A nonnegative, unit-independent electric-current magnitude
\end{definition}
\begin{proof}
  This is the declared model definition of Electric Current Quantity.
\end{proof}
\begin{definition}[Magnetic Permeability Quantity]
  \label{def:physics:phyx-mini-0916:phyxminiproblems-problemphyxmini0916-magneticpermeabilityquantity}
  \lean{PhyXMiniProblems.ProblemPhyXMini0916.MagneticPermeabilityQuantity}
  \uses{def:physics:phyx-mini-0916:phyxminiproblems-problemphyxmini0916-magneticpermeabilitydimension}
  A nonnegative, unit-independent magnetic permeability
\end{definition}
\begin{proof}
  This is the declared model definition of Magnetic Permeability Quantity.
\end{proof}
\begin{definition}[length In Meters]
  \label{def:physics:phyx-mini-0916:phyxminiproblems-problemphyxmini0916-lengthinmeters}
  \lean{PhyXMiniProblems.ProblemPhyXMini0916.lengthInMeters}
  \uses{def:physics:phyx-mini-0916:phyxminiproblems-problemphyxmini0916-lengthquantity}
  SI metre readout of a physical length
\end{definition}
\begin{proof}
  This is the declared model definition of length In Meters.
\end{proof}
\begin{definition}[length In Centimeters]
  \label{def:physics:phyx-mini-0916:phyxminiproblems-problemphyxmini0916-lengthincentimeters}
  \lean{PhyXMiniProblems.ProblemPhyXMini0916.lengthInCentimeters}
  \uses{def:physics:phyx-mini-0916:phyxminiproblems-problemphyxmini0916-lengthquantity, def:physics:phyx-mini-0916:phyxminiproblems-problemphyxmini0916-lengthinmeters}
  Centimetre readout used by the printed segment-length label
\end{definition}
\begin{proof}
  This is the declared model definition of length In Centimeters.
\end{proof}
\begin{definition}[current In Amperes]
  \label{def:physics:phyx-mini-0916:phyxminiproblems-problemphyxmini0916-currentinamperes}
  \lean{PhyXMiniProblems.ProblemPhyXMini0916.currentInAmperes}
  \uses{def:physics:phyx-mini-0916:phyxminiproblems-problemphyxmini0916-electriccurrentquantity}
  SI ampere readout of a physical electric-current magnitude
\end{definition}
\begin{proof}
  This is the declared model definition of current In Amperes.
\end{proof}
\begin{definition}[permeability In Tesla Meters Per Ampere]
  \label{def:physics:phyx-mini-0916:phyxminiproblems-problemphyxmini0916-permeabilityinteslametersperampere}
  \lean{PhyXMiniProblems.ProblemPhyXMini0916.permeabilityInTeslaMetersPerAmpere}
  \uses{def:physics:phyx-mini-0916:phyxminiproblems-problemphyxmini0916-magneticpermeabilityquantity}
  SI readout of permeability in tesla-metres per ampere
\end{definition}
\begin{proof}
  This is the declared model definition of permeability In Tesla Meters Per Ampere.
\end{proof}
\begin{definition}[Spatial Vector]
  \label{def:physics:phyx-mini-0916:phyxminiproblems-problemphyxmini0916-spatialvector}
  \lean{PhyXMiniProblems.ProblemPhyXMini0916.SpatialVector}
  A three-dimensional spatial vector
\end{definition}
\begin{proof}
  This is the declared model definition of Spatial Vector.
\end{proof}
\begin{definition}[spatial Cross Product]
  \label{def:physics:phyx-mini-0916:phyxminiproblems-problemphyxmini0916-spatialcrossproduct}
  \lean{PhyXMiniProblems.ProblemPhyXMini0916.spatialCrossProduct}
  \uses{def:physics:phyx-mini-0916:phyxminiproblems-problemphyxmini0916-spatialvector}
  Mathlib's coordinate cross product transported to the Euclidean-space type used by Physlib magnetic fields
\end{definition}
\begin{proof}
  This is the declared model definition of spatial Cross Product.
\end{proof}
\begin{definition}[Coordinate Axis]
  \label{def:physics:phyx-mini-0916:phyxminiproblems-problemphyxmini0916-coordinateaxis}
  \lean{PhyXMiniProblems.ProblemPhyXMini0916.CoordinateAxis}
  The three coordinate axes explicitly labelled in the image
\end{definition}
\begin{proof}
  This is the declared model definition of Coordinate Axis.
\end{proof}
\begin{definition}[axis Direction]
  \label{def:physics:phyx-mini-0916:phyxminiproblems-problemphyxmini0916-axisdirection}
  \lean{PhyXMiniProblems.ProblemPhyXMini0916.axisDirection}
  \uses{def:physics:phyx-mini-0916:phyxminiproblems-problemphyxmini0916-spatialvector, def:physics:phyx-mini-0916:phyxminiproblems-problemphyxmini0916-coordinateaxis}
  Unit vector associated with a displayed coordinate axis
\end{definition}
\begin{proof}
  This is the declared model definition of axis Direction.
\end{proof}
\begin{definition}[Observation Point]
  \label{def:physics:phyx-mini-0916:phyxminiproblems-problemphyxmini0916-observationpoint}
  \lean{PhyXMiniProblems.ProblemPhyXMini0916.ObservationPoint}
  The two observation points named in the primary figure
\end{definition}
\begin{proof}
  This is the declared model definition of Observation Point.
\end{proof}
\begin{definition}[Current Arrow Site]
  \label{def:physics:phyx-mini-0916:phyxminiproblems-problemphyxmini0916-currentarrowsite}
  \lean{PhyXMiniProblems.ProblemPhyXMini0916.CurrentArrowSite}
  The two pictured locations of the repeated current arrow
\end{definition}
\begin{proof}
  This is the declared model definition of Current Arrow Site.
\end{proof}
\begin{definition}[Horizontal Direction]
  \label{def:physics:phyx-mini-0916:phyxminiproblems-problemphyxmini0916-horizontaldirection}
  \lean{PhyXMiniProblems.ProblemPhyXMini0916.HorizontalDirection}
  Direction of the current arrows along the horizontal wire
\end{definition}
\begin{proof}
  This is the declared model definition of Horizontal Direction.
\end{proof}
\begin{definition}[Wire Material]
  \label{def:physics:phyx-mini-0916:phyxminiproblems-problemphyxmini0916-wirematerial}
  \lean{PhyXMiniProblems.ProblemPhyXMini0916.WireMaterial}
  Material of the current-carrying wire
\end{definition}
\begin{proof}
  This is the declared model definition of Wire Material.
\end{proof}
\begin{definition}[Current Destination]
  \label{def:physics:phyx-mini-0916:phyxminiproblems-problemphyxmini0916-currentdestination}
  \lean{PhyXMiniProblems.ProblemPhyXMini0916.CurrentDestination}
  Destination stated for the current in the problem prose
\end{definition}
\begin{proof}
  This is the declared model definition of Current Destination.
\end{proof}
\begin{definition}[Current Regime]
  \label{def:physics:phyx-mini-0916:phyxminiproblems-problemphyxmini0916-currentregime}
  \lean{PhyXMiniProblems.ProblemPhyXMini0916.CurrentRegime}
  Whether the current is time-independent
\end{definition}
\begin{proof}
  This is the declared model definition of Current Regime.
\end{proof}
\begin{definition}[Short Wire Figure]
  \label{def:physics:phyx-mini-0916:phyxminiproblems-problemphyxmini0916-shortwirefigure}
  \lean{PhyXMiniProblems.ProblemPhyXMini0916.ShortWireFigure}
  \uses{def:physics:phyx-mini-0916:phyxminiproblems-problemphyxmini0916-coordinateaxis, def:physics:phyx-mini-0916:phyxminiproblems-problemphyxmini0916-observationpoint, def:physics:phyx-mini-0916:phyxminiproblems-problemphyxmini0916-currentarrowsite, def:physics:phyx-mini-0916:phyxminiproblems-problemphyxmini0916-horizontaldirection}
  Literal labels and qualitative objects in the supplied image 916.png. The displayed scalar labels are kept separate from the physical quantities that they calibrate
\end{definition}
\begin{proof}
  This is the declared model definition of Short Wire Figure.
\end{proof}
\begin{definition}[Short Wire Magnetic Field Setup]
  \label{def:physics:phyx-mini-0916:phyxminiproblems-problemphyxmini0916-shortwiremagneticfieldsetup}
  \lean{PhyXMiniProblems.ProblemPhyXMini0916.ShortWireMagneticFieldSetup}
  \uses{def:physics:phyx-mini-0916:phyxminiproblems-problemphyxmini0916-lengthquantity, def:physics:phyx-mini-0916:phyxminiproblems-problemphyxmini0916-electriccurrentquantity, def:physics:phyx-mini-0916:phyxminiproblems-problemphyxmini0916-magneticpermeabilityquantity, def:physics:phyx-mini-0916:phyxminiproblems-problemphyxmini0916-spatialvector, def:physics:phyx-mini-0916:phyxminiproblems-problemphyxmini0916-observationpoint, def:physics:phyx-mini-0916:phyxminiproblems-problemphyxmini0916-wirematerial, def:physics:phyx-mini-0916:phyxminiproblems-problemphyxmini0916-currentdestination, def:physics:phyx-mini-0916:phyxminiproblems-problemphyxmini0916-currentregime, def:physics:phyx-mini-0916:phyxminiproblems-problemphyxmini0916-shortwirefigure}
  Independent physical quantities and observables for the wire segment. The field is not defined from an answer choice or from the requested numerical component
\end{definition}
\begin{proof}
  This is the declared model definition of Short Wire Magnetic Field Setup.
\end{proof}
\begin{definition}[magnetic Field Vector At P2 In Teslas]
  \label{def:physics:phyx-mini-0916:phyxminiproblems-problemphyxmini0916-magneticfieldvectoratp2inteslas}
  \lean{PhyXMiniProblems.ProblemPhyXMini0916.magneticFieldVectorAtP2InTeslas}
  \uses{def:physics:phyx-mini-0916:phyxminiproblems-problemphyxmini0916-spatialvector, def:physics:phyx-mini-0916:phyxminiproblems-problemphyxmini0916-shortwiremagneticfieldsetup}
  The Physlib magnetic-field vector at P₂, read in coherent SI units
\end{definition}
\begin{proof}
  This is the declared model definition of magnetic Field Vector At P2 In Teslas.
\end{proof}
\begin{definition}[magnetic Field ZComponent At P2 In Teslas]
  \label{def:physics:phyx-mini-0916:phyxminiproblems-problemphyxmini0916-magneticfieldzcomponentatp2inteslas}
  \lean{PhyXMiniProblems.ProblemPhyXMini0916.magneticFieldZComponentAtP2InTeslas}
  \uses{def:physics:phyx-mini-0916:phyxminiproblems-problemphyxmini0916-shortwiremagneticfieldsetup, def:physics:phyx-mini-0916:phyxminiproblems-problemphyxmini0916-magneticfieldvectoratp2inteslas}
  Signed z component of the field at P₂, in teslas
\end{definition}
\begin{proof}
  This is the declared model definition of magnetic Field ZComponent At P2 In Teslas.
\end{proof}
\begin{definition}[Matches Short Wire Problem Data]
  \label{def:physics:phyx-mini-0916:phyxminiproblems-problemphyxmini0916-matchesshortwireproblemdata}
  \lean{PhyXMiniProblems.ProblemPhyXMini0916.MatchesShortWireProblemData}
  \uses{def:physics:phyx-mini-0916:phyxminiproblems-problemphyxmini0916-lengthinmeters, def:physics:phyx-mini-0916:phyxminiproblems-problemphyxmini0916-lengthincentimeters, def:physics:phyx-mini-0916:phyxminiproblems-problemphyxmini0916-currentinamperes, def:physics:phyx-mini-0916:phyxminiproblems-problemphyxmini0916-shortwiremagneticfieldsetup}
  Prose data and the physical quantities calibrated by its numerical values
\end{definition}
\begin{proof}
  This is the declared model definition of Matches Short Wire Problem Data.
\end{proof}
\begin{definition}[Matches Primary Short Wire Figure]
  \label{def:physics:phyx-mini-0916:phyxminiproblems-problemphyxmini0916-matchesprimaryshortwirefigure}
  \lean{PhyXMiniProblems.ProblemPhyXMini0916.MatchesPrimaryShortWireFigure}
  \uses{def:physics:phyx-mini-0916:phyxminiproblems-problemphyxmini0916-lengthinmeters, def:physics:phyx-mini-0916:phyxminiproblems-problemphyxmini0916-lengthincentimeters, def:physics:phyx-mini-0916:phyxminiproblems-problemphyxmini0916-currentinamperes, def:physics:phyx-mini-0916:phyxminiproblems-problemphyxmini0916-axisdirection, def:physics:phyx-mini-0916:phyxminiproblems-problemphyxmini0916-shortwiremagneticfieldsetup}
  Primary-image transcription and the corresponding metre-coordinate geometry. The ray to P₂ is 30° above the negative x direction. Combined with the leftward current this records the sign information needed by the right-hand rule, but it does not assert any field value
\end{definition}
\begin{proof}
  This is the declared model definition of Matches Primary Short Wire Figure.
\end{proof}
\begin{definition}[Has Physical Short Wire Parameters]
  \label{def:physics:phyx-mini-0916:phyxminiproblems-problemphyxmini0916-hasphysicalshortwireparameters}
  \lean{PhyXMiniProblems.ProblemPhyXMini0916.HasPhysicalShortWireParameters}
  \uses{def:physics:phyx-mini-0916:phyxminiproblems-problemphyxmini0916-lengthinmeters, def:physics:phyx-mini-0916:phyxminiproblems-problemphyxmini0916-currentinamperes, def:physics:phyx-mini-0916:phyxminiproblems-problemphyxmini0916-permeabilityinteslametersperampere, def:physics:phyx-mini-0916:phyxminiproblems-problemphyxmini0916-shortwiremagneticfieldsetup}
  Positivity and unit-vector assumptions selecting the physical branch
\end{definition}
\begin{proof}
  This is the declared model definition of Has Physical Short Wire Parameters.
\end{proof}
\begin{definition}[Uses Standard Vacuum Permeability]
  \label{def:physics:phyx-mini-0916:phyxminiproblems-problemphyxmini0916-usesstandardvacuumpermeability}
  \lean{PhyXMiniProblems.ProblemPhyXMini0916.UsesStandardVacuumPermeability}
  \uses{def:physics:phyx-mini-0916:phyxminiproblems-problemphyxmini0916-permeabilityinteslametersperampere, def:physics:phyx-mini-0916:phyxminiproblems-problemphyxmini0916-shortwiremagneticfieldsetup}
  The dimensionful permeability agrees with the Physlib electromagnetic system and has the standard vacuum SI value 4π * 10⁻⁷ T m/A. This is calibration data, not the requested magnetic field
\end{definition}
\begin{proof}
  This is the declared model definition of Uses Standard Vacuum Permeability.
\end{proof}
\begin{definition}[Satisfies Steady Current Field Model]
  \label{def:physics:phyx-mini-0916:phyxminiproblems-problemphyxmini0916-satisfiessteadycurrentfieldmodel}
  \lean{PhyXMiniProblems.ProblemPhyXMini0916.SatisfiesSteadyCurrentFieldModel}
  \uses{def:physics:phyx-mini-0916:phyxminiproblems-problemphyxmini0916-shortwiremagneticfieldsetup}
  A steady current produces a time-independent field in this model
\end{definition}
\begin{proof}
  This is the declared model definition of Satisfies Steady Current Field Model.
\end{proof}
\begin{definition}[Satisfies Short Segment Biot Savart Law]
  \label{def:physics:phyx-mini-0916:phyxminiproblems-problemphyxmini0916-satisfiesshortsegmentbiotsavartlaw}
  \lean{PhyXMiniProblems.ProblemPhyXMini0916.SatisfiesShortSegmentBiotSavartLaw}
  \uses{def:physics:phyx-mini-0916:phyxminiproblems-problemphyxmini0916-lengthinmeters, def:physics:phyx-mini-0916:phyxminiproblems-problemphyxmini0916-currentinamperes, def:physics:phyx-mini-0916:phyxminiproblems-problemphyxmini0916-permeabilityinteslametersperampere, def:physics:phyx-mini-0916:phyxminiproblems-problemphyxmini0916-spatialcrossproduct, def:physics:phyx-mini-0916:phyxminiproblems-problemphyxmini0916-shortwiremagneticfieldsetup, def:physics:phyx-mini-0916:phyxminiproblems-problemphyxmini0916-magneticfieldvectoratp2inteslas}
  Short-current-element form of the Biot--Savart law at P₂: B = μ₀ I Δℓ /(4πr²) (î\_current × r̂). The cross product retains both the sin θ factor and the right-hand-rule sign. This premise is a general governing relation in the stored physical quantities; it contains neither the requested numerical field nor an answer choice
\end{definition}
\begin{proof}
  This is the declared model definition of Satisfies Short Segment Biot Savart Law.
\end{proof}
\begin{lemma}[magnetic Field Vector At P2 exact]
  \label{lem:physics:phyx-mini-0916:phyxminiproblems-problemphyxmini0916-magneticfieldvectoratp2-exact}
  \lean{PhyXMiniProblems.ProblemPhyXMini0916.magneticFieldVectorAtP2_exact}
  \uses{def:physics:phyx-mini-0916:phyxminiproblems-problemphyxmini0916-axisdirection, def:physics:phyx-mini-0916:phyxminiproblems-problemphyxmini0916-shortwiremagneticfieldsetup, def:physics:phyx-mini-0916:phyxminiproblems-problemphyxmini0916-magneticfieldvectoratp2inteslas, def:physics:phyx-mini-0916:phyxminiproblems-problemphyxmini0916-matchesshortwireproblemdata, def:physics:phyx-mini-0916:phyxminiproblems-problemphyxmini0916-matchesprimaryshortwirefigure, def:physics:phyx-mini-0916:phyxminiproblems-problemphyxmini0916-hasphysicalshortwireparameters, def:physics:phyx-mini-0916:phyxminiproblems-problemphyxmini0916-usesstandardvacuumpermeability, def:physics:phyx-mini-0916:phyxminiproblems-problemphyxmini0916-satisfiessteadycurrentfieldmodel, def:physics:phyx-mini-0916:phyxminiproblems-problemphyxmini0916-satisfiesshortsegmentbiotsavartlaw}
  The short-segment Biot--Savart calculation gives the exact model value -(1 / 23040000) T in the z direction. Its decimal magnitude is about 4.3403 * 10⁻⁸ T
\end{lemma}
\begin{proof}
  The conclusion follows from the governing relations and figure readouts named in the statement of magnetic Field Vector At P2 exact.
\end{proof}
\begin{definition}[Answer Choice]
  \label{def:physics:phyx-mini-0916:phyxminiproblems-problemphyxmini0916-answerchoice}
  \lean{PhyXMiniProblems.ProblemPhyXMini0916.AnswerChoice}
  Labels of the four signed magnetic-field answers in the question
\end{definition}
\begin{proof}
  This is the declared model definition of Answer Choice.
\end{proof}
\begin{definition}[displayed Magnetic Field In Teslas]
  \label{def:physics:phyx-mini-0916:phyxminiproblems-problemphyxmini0916-displayedmagneticfieldinteslas}
  \lean{PhyXMiniProblems.ProblemPhyXMini0916.displayedMagneticFieldInTeslas}
  \uses{def:physics:phyx-mini-0916:phyxminiproblems-problemphyxmini0916-answerchoice}
  Signed magnetic-field component in teslas displayed beside each choice
\end{definition}
\begin{proof}
  This is the declared model definition of displayed Magnetic Field In Teslas.
\end{proof}
\begin{definition}[Rounds To Nearest Tenth Times Ten Pow Neg Eight]
  \label{def:physics:phyx-mini-0916:phyxminiproblems-problemphyxmini0916-roundstonearesttenthtimestenpownegeight}
  \lean{PhyXMiniProblems.ProblemPhyXMini0916.RoundsToNearestTenthTimesTenPowNegEight}
  Rounding to the nearest 0.1 * 10⁻⁸ T, the precision used by choices A, C, and D
\end{definition}
\begin{proof}
  This is the declared model definition of Rounds To Nearest Tenth Times Ten Pow Neg Eight.
\end{proof}
\begin{definition}[Is Unique Closest Displayed Answer]
  \label{def:physics:phyx-mini-0916:phyxminiproblems-problemphyxmini0916-isuniqueclosestdisplayedanswer}
  \lean{PhyXMiniProblems.ProblemPhyXMini0916.IsUniqueClosestDisplayedAnswer}
  \uses{def:physics:phyx-mini-0916:phyxminiproblems-problemphyxmini0916-answerchoice, def:physics:phyx-mini-0916:phyxminiproblems-problemphyxmini0916-displayedmagneticfieldinteslas}
  A choice is strictly closer than every distinct displayed alternative
\end{definition}
\begin{proof}
  This is the declared model definition of Is Unique Closest Displayed Answer.
\end{proof}
% --- Archon declaration topology end ---
% --- Archon physics formalization source end ---
